\subsection{Dedekind 整环的刻画}

定理1. 设 A 为整环，K 为其分式域。以下条件等价：

\begin{enumerate}
\def\labelenumi{(\alph{enumi})}
\item
  A 是 Dedekind 整环；
\item
  A 是 Krull 整环，且 K 上每个在 A 上取正值的非平凡赋值都等价于 A 的一个本性赋值；
\item
  A 是 Krull 整环，且 A 的每个分式理想 \(\Im \neq (0)\) 都是除性理想；
\item
  A 的每个分式理想 \(\Im \neq (0)\) 都是可逆的；
\item
  A 是诺特整闭环，且 A 的每个非零素理想都是极大理想；
\item
  A 是诺特环，并且对 A 的每个极大理想 m，\(A_{m}\) 或为域或为离散赋值环；
\item
  A 是诺特环，并且对 A 的每个极大理想 m，\(A_{m}\) 是主理想整环。
\end{enumerate}

我们先证明 (a) 与 (b) 等价。第1节第6号定理3的推论2立即表明 (a) 蕴含 (b)。反过来，(b) 蕴含 (a)，因为对于 A 的每个素理想 p，都存在支配 \(A_{p}\) 的 K 的赋值环（第VI章，第1节，第2号，定理2的推论）。

证明的其余部分通过证明以下蕴含来完成：

\[
(a) \Rightarrow (c) \Rightarrow (d) \Rightarrow (e) \Rightarrow (f) \Rightarrow (g) \Rightarrow (a).
\]

若 A 是 Dedekind 整环且 b 是非零分式理想，则对于每个极大理想 p，有 \(bA_{p} = \tilde{b}A_{p}\)（第1节，第4号，命题7），从而 \(b = \tilde{b}\)（第II章，第3节，第3号，定理1的推论3）；因此 (a) 蕴含 (c)。现在证明 (c) 蕴含 (d)。若 (c) 成立，映射 \(a \mapsto div a\) 是从 I(A) 到 D(A) 的双射（参见第1节，第1号）；它是一个同态（第1节，第2号），而 D(A) 是一个群，所以 I(A) 的每个元素都可逆。

我们证明 (d) 蕴含 (e)。若 (d) 成立，A 的每个整理想（\(\neq(0)\)）都是有限生成的（第II章，第5节，第6号，定理4），因而 A 是诺特环；由于 I(A) 是一个群，D(A) 也是一个群，所以 A 是完全整闭的（第1节，第2号，定理1）。最后，若 p 是 A 的非零素理想，且 m 是包含 p 的 A 的极大理想，则 \(A_{m}\) 是主理想整环（第II章，第5节，第6号，定理4）；由于 \(pA_{m}\) 是非零素理想，必有 \(pA_{m}=mA_{m}\)（主理想整环是 Dedekind 整环），从而 p=m（第II章，第2节，第5号，命题11），p 也是极大理想。

现在证明 (e) 蕴含 (f)。若 m 是 A 的极大理想且 (e) 成立，则 \(A_{m}\) 是整闭的诺特环，而其极大理想 \(mA_{m}\) 或为 (0)，或是 \(A_{m}\) 的唯一非零素理想；因此根据第1节第7号的命题11，\(A_{m}\) 是域或离散赋值环。

显然 (f) 蕴含 (g)。

最后证明 (g) 蕴含 (a)。由于 A 是所有 \(A_{m}\) 的交，其中 m 遍历极大理想集合（第II章，第3节，第3号，定理1的推论4），(g) 蕴含 A 是整闭的且为诺特环，因而 A 是 Krull 整环（第1节，第3号，定理2的推论）。另一方面，可以像证明 (d) \(\Rightarrow\) (e) 时那样证明 A 的每个非零素理想都是极大理想。

命题1. 半局部 Dedekind 整环是主理想整环。

设 A 为半局部 Dedekind 整环，K 为其分式域，\(p_{1}, \ldots, p_{n}\) 为其极大理想，\(v_{1}, \ldots, v_{n}\) 为相应的本性赋值；这些是 A 的唯一本性赋值。设 a 为 A 的非零整理想。由于它是除性理想，存在（第1节，第4号，命题5）整数 \(q_{1}, \ldots, q_{n}\)，使 a 是满足 \(x \in K\) 且 \(v_{i}(x) \geqslant q_{i}\)（\(1 \leqslant i \leqslant n\)）的元素集合。设 \(x_{0}\) 是 K 中满足 \(v_{i}(x_{0}) = q_{i}\)（\(1 \leqslant i \leqslant n\)）的元素（第VI章，第7节，第2号，定理1的推论1）。于是 a 是满足 \(x \in K\) 且 \(v_{i}(xx_{0}^{-1}) \geqslant 0\)（\(1 \leqslant i \leqslant n\)）的元素集合。因此 \(a = Ax_{0}\)。

如果 A 是 Dedekind 整环，那么在定理1的证明中已经看到，A 的除子群 D(A) 可与分式理想群 I(A) 等同，其中 \(a \neq (0)\)（因为 A 是诺特环，每个非零分式理想都是有限生成的）。A 的除子类群 C(A)（\(\S 1\)，第2号）于是可与 A 的非零理想类群 \(\neq 0\) 等同（该类群定义于第II章，\(\S 5\)，第7号）。

